\documentclass[12pt,a4paper]{article}
\usepackage[utf8]{inputenc}
\usepackage[T5]{fontenc}
\usepackage[vietnamese]{babel}
\usepackage[margin=1.3cm,top=0.8cm,bottom=1.0cm,includehead,includefoot,headheight=18pt,headsep=6pt,footskip=16pt]{geometry}
\usepackage{amsmath,amssymb,array,tabularx,enumitem,fancyhdr,tikz,xcolor,needspace,multicol}
\usetikzlibrary{calc,angles,quotes,arrows.meta,babel}

\definecolor{ink}{RGB}{30,55,70}
\definecolor{accent}{RGB}{18,82,112}

\setlength{\parindent}{0pt}
\setlength{\parskip}{2pt}

\newcommand{\headerQuote}{%
  \ifcase\numexpr\value{page}-1\relax
    Mathematics is about understanding, not just numbers%
  \or The essence of mathematics lies in its freedom%
  \or Pure mathematics is the poetry of logical ideas%
  \or In mathematics, asking good questions matters most%
  \or Everything around you is built on mathematics%
  \or Mathematics is the queen of the sciences%
  \or Stay with tough problems a little bit longer%
  \or Math gives us hope that every problem has a solution%
  \or Success is the sum of small daily efforts%
  \or Think deeply, check carefully, solve cleanly%
  \or Genius is 1\% inspiration and 99\% perspiration%
  \or Practice makes perfect — keep solving!%
  \or Mathematics is about understanding, not just numbers%
  \or The essence of mathematics lies in its freedom%
  \or Pure mathematics is the poetry of logical ideas%
  \or In mathematics, asking good questions matters most%
  \or Everything around you is built on mathematics%
  \or Mathematics is the queen of the sciences%
  \or Stay with tough problems a little bit longer%
  \or Math gives us hope that every problem has a solution%
  \or Success is the sum of small daily efforts%
  \or Think deeply, check carefully, solve cleanly%
  \or Genius is 1\% inspiration and 99\% perspiration%
  \or Practice makes perfect — keep solving!%
  \or Mathematics is about understanding, not just numbers%
  \or The essence of mathematics lies in its freedom%
  \or Pure mathematics is the poetry of logical ideas%
  \or In mathematics, asking good questions matters most%
  \or Everything around you is built on mathematics%
  \or Mathematics is the queen of the sciences%
  \or Stay with tough problems a little bit longer%
  \or Math gives us hope that every problem has a solution%
  \or Success is the sum of small daily efforts%
  \or Think deeply, check carefully, solve cleanly%
  \or Genius is 1\% inspiration and 99\% perspiration%
  \or Practice makes perfect — keep solving!%
  \or Mathematics is about understanding, not just numbers%
  \or The essence of mathematics lies in its freedom%
  \or Pure mathematics is the poetry of logical ideas%
  \or In mathematics, asking good questions matters most%
  \else Practice makes perfect — keep solving!%
  \fi
}

\pagestyle{fancy}
\fancyhf{}
\lhead{\footnotesize\textit{\color{gray!70!black}\headerQuote}}
\rhead{\small\textbf{\color{ink}Người biên soạn: Dương Gia Bảo}}
\lfoot{}
\rfoot{\small Trang \thepage/12}
\renewcommand{\headrulewidth}{0.4pt}
\renewcommand{\footrulewidth}{0pt}

\fancypagestyle{dapan}{
  \fancyhf{}
  \lhead{\footnotesize\textit{\color{gray!70!black}\headerQuote}}
  \rhead{\small\textbf{\color{ink}Người biên soạn: Dương Gia Bảo}}
  \lfoot{\small\textit{\color{gray!60!black}Đáp án Đề cương Ôn tập Chương VI (50 câu) — Toán 10 KNTT}}
  \rfoot{\small Trang \thepage/2}
  \renewcommand{\headrulewidth}{0.4pt}
  \renewcommand{\footrulewidth}{0pt}
}

\newcommand{\chapterbanner}[2]{%
  \par\vspace{3pt}\noindent
  \begin{tikzpicture}
    \node[draw=accent, fill=accent!6, rounded corners=4pt, inner xsep=10pt, inner ysep=6pt, text width=\dimexpr\linewidth-20.8pt\relax, line width=0.8pt] {%
      {\large\bfseries\color{ink}#1}\\[2pt]
      {\small\color{gray!75!black}#2}%
    };
  \end{tikzpicture}\par\nopagebreak\vspace{5pt}
}

\newcommand{\parttitle}[1]{\par\vspace{3pt}\needspace{3\baselineskip}{\normalsize\bfseries\color{accent}#1}\par\nopagebreak\vspace{2pt}}
\newcommand{\four}[4]{\par\nopagebreak\vspace{1.5pt}\noindent\begin{tabularx}{\linewidth}{@{}*{4}{>{\raggedright\arraybackslash}X}@{}}\textbf{A.} #1&\textbf{B.} #2&\textbf{C.} #3&\textbf{D.} #4\end{tabularx}}
\newcommand{\two}[4]{\par\nopagebreak\vspace{1.5pt}\noindent\begin{tabularx}{\linewidth}{@{}*{2}{>{\raggedright\arraybackslash}X}@{}}\textbf{A.} #1&\textbf{B.} #2\\[3.8pt]\textbf{C.} #3&\textbf{D.} #4\end{tabularx}}

\newcommand{\answerbox}{%
  {\unskip\nobreak\hfil\penalty50\hskip1.2em\hbox{}\nobreak\hfil
  \mbox{\textit{Đáp số:}\ \fbox{\rule{0pt}{8.5pt}\hspace{2.2cm}}}%
  \parfillskip=0pt\finalhyphendemerits=0\par}%
}

\setlist[enumerate,1]{label=\textbf{Câu \arabic*.},leftmargin=1.15cm,labelsep=.15cm,itemsep=2pt,topsep=2pt,parsep=0pt}

\newcommand{\oli}[1]{%
  \par\nopagebreak\vspace{2.5pt}\noindent
  \begin{tikzpicture}
    \draw[step=0.5cm, color=cyan!25!gray!25, line width=0.25pt] (0,0) grid (\linewidth, #1);
    \draw[color=cyan!50!gray!45, line width=0.45pt, rounded corners=2pt] (0,0) rectangle (\linewidth, #1);
    \node[anchor=north west, font=\scriptsize\itshape\color{gray!75}, inner sep=3pt] at (0, #1) {Bài làm:};
  \end{tikzpicture}\par\vspace{3.5pt}
}

\newcolumntype{C}[1]{>{\centering\arraybackslash}m{#1}}


\begin{document}

\chapterbanner{ĐỀ CƯƠNG ÔN TẬP CHƯƠNG VI: HÀM SỐ, ĐỒ THỊ VÀ ỨNG DỤNG}{Môn Toán 10 — Kết nối tri thức với cuộc sống \quad|\quad Tổng cộng: 50 câu (20 Trắc nghiệm + 30 Trả lời ngắn)}

\parttitle{PHẦN I. CÂU TRẮC NGHIỆM NHIỀU PHƯƠNG ÁN LỰA CHỌN (20 CÂU)}
\textit{Thí sinh khoanh tròn vào chữ cái đứng trước phương án trả lời đúng nhất.}

\parttitle{Bài 15. Hàm số và đồ thị}

\begin{enumerate}

\needspace{4\baselineskip}\item Tập xác định $D$ của hàm số $y = \dfrac{\sqrt{x + 2}}{x^2 - 1}$ là
\two{$(-2; +\infty) \setminus \{1\}$.}{$[-2; +\infty) \setminus \{1\}$.}{$[-2; +\infty) \setminus \{-1; 1\}$.}{$\mathbb{R} \setminus \{-1; 1\}$.}

\needspace{4\baselineskip}\item Điểm nào sau đây thuộc đồ thị hàm số $y = \dfrac{2x - 1}{\sqrt{x + 3}}$?
\four{$M\left(1; \dfrac{1}{2}\right)$.}{$N(-3; 0)$.}{$P(6; 3)$.}{$Q(1; 1)$.}

\needspace{4\baselineskip}\item Cho hàm số $f(x) = \begin{cases} 3x - 2 & \text{khi } x \ge 2 \\ x^2 + 1 & \text{khi } x < 2 \end{cases}$. Giá trị của biểu thức $P = f(4) + f(-1)$ bằng
\four{$10$.}{$8$.}{$14$.}{$12$.}

\needspace{4\baselineskip}\item Cho hàm số $y = f(x)$ có đồ thị là đường gấp khúc đi qua các điểm $A(-3; 1)$, $B(-1; 4)$, $C(2; 0)$, $D(5; 3)$. Hàm số nghịch biến trên khoảng nào sau đây?
\four{$(-3; -1)$.}{$(-1; 2)$.}{$(2; 5)$.}{$(-3; 2)$.}

\needspace{4\baselineskip}\item Tập giá trị của hàm số $y = 4 - \sqrt{9 - x^2}$ là
\four{$[0; 3]$.}{$[0; 4]$.}{$(-\infty; 4]$.}{$[1; 4]$.}

\end{enumerate}

\parttitle{Bài 16. Hàm số bậc hai}

\begin{enumerate}[resume*]

\needspace{4\baselineskip}\item Tọa độ đỉnh $I$ của parabol $(P): y = -2x^2 + 8x - 3$ là
\four{$I(-2; -27)$.}{$I(2; 5)$.}{$I(4; -3)$.}{$I(2; 13)$.}

\needspace{4\baselineskip}\item Hàm số bậc hai $y = 3x^2 - 6x + 4$ đồng biến trên khoảng nào sau đây?
\four{$(1; +\infty)$.}{$(-\infty; 1)$.}{$(-1; +\infty)$.}{$(-\infty; 2)$.}

\needspace{4\baselineskip}\item Parabol $(P): y = x^2 - 5x + 6$ cắt trục hoành tại hai điểm phân biệt $A, B$. Độ dài đoạn thẳng $AB$ bằng
\four{$5$.}{$6$.}{$1$.}{$2$.}

\needspace{4\baselineskip}\item Cho parabol $(P): y = ax^2 + bx + c$ ($a \ne 0$) quay bề lõm lên trên, có đỉnh $I(-1; -4)$ và cắt trục tung tại điểm có tung độ bằng $-3$. Khẳng định nào sau đây đúng?
\two{$a > 0,\; b < 0,\; c < 0$.}{$a > 0,\; b > 0,\; c < 0$.}{$a < 0,\; b > 0,\; c < 0$.}{$a > 0,\; b > 0,\; c > 0$.}

\needspace{4\baselineskip}\item Giá trị lớn nhất của hàm số $y = x^2 - 2x - 3$ trên đoạn $[0; 4]$ bằng
\four{$-4$.}{$-3$.}{$1$.}{$5$.}

\end{enumerate}

\newpage

\begin{enumerate}[resume*]

\needspace{4\baselineskip}\item Biết parabol $(P): y = x^2 + bx + c$ có trục đối xứng $x = 3$ và đi qua điểm $M(1; 0)$. Giá trị của tổng $b + c$ bằng
\four{$-1$.}{$1$.}{$-11$.}{$11$.}

\end{enumerate}

\parttitle{Bài 17. Dấu của tam thức bậc hai}

\begin{enumerate}[resume*]

\needspace{4\baselineskip}\item Tam thức bậc hai nào sau đây nhận giá trị dương với mọi $x \in \mathbb{R}$?
\two{$f(x) = x^2 - 4x + 4$.}{$f(x) = -x^2 + 4x - 5$.}{$f(x) = x^2 - 4x + 5$.}{$f(x) = x^2 - 5x + 4$.}

\needspace{4\baselineskip}\item Cho tam thức bậc hai $f(x) = -2x^2 + 5x - 2$. Tập hợp tất cả các giá trị của $x$ để $f(x) > 0$ là
\two{$\left(-\infty; \dfrac{1}{2}\right) \cup (2; +\infty)$.}{$\left[\dfrac{1}{2}; 2\right]$.}{$(-\infty; 2)$.}{$\left(\dfrac{1}{2}; 2\right)$.}

\needspace{4\baselineskip}\item Tập nghiệm $S$ của bất phương trình $3x^2 - 7x + 2 \ge 0$ là
\two{$\left(-\infty; \dfrac{1}{3}\right] \cup [2; +\infty)$.}{$\left[\dfrac{1}{3}; 2\right]$.}{$\left(-\infty; \dfrac{1}{3}\right) \cup (2; +\infty)$.}{$\left(\dfrac{1}{3}; 2\right)$.}

\needspace{4\baselineskip}\item Số nghiệm nguyên của bất phương trình $x^2 - 4x - 5 < 0$ là
\four{$7$.}{$6$.}{$5$.}{$4$.}

\needspace{4\baselineskip}\item Tất cả các giá trị của tham số $m$ để tam thức $f(x) = x^2 - 2mx + m + 2$ luôn dương với mọi $x \in \mathbb{R}$ là
\two{$m \in [-1; 2]$.}{$m \in (-1; 2)$.}{$m \in (-\infty; -1) \cup (2; +\infty)$.}{$m \in (-2; 1)$.}

\end{enumerate}

\parttitle{Bài 18. Phương trình quy về phương trình bậc hai}

\begin{enumerate}[resume*]

\needspace{4\baselineskip}\item Tập nghiệm của phương trình $\sqrt{x^2 + 3x - 2} = \sqrt{1 + x}$ là
\four{$\{1\}$.}{$\{-3\}$.}{$\{-3; 1\}$.}{$\varnothing$.}

\needspace{4\baselineskip}\item Số nghiệm thực của phương trình $\sqrt{2x^2 - 4x + 9} = x + 1$ là
\four{$0$.}{$1$.}{$3$.}{$2$.}

\needspace{4\baselineskip}\item Tổng các nghiệm thực của phương trình $\sqrt{x^2 - 4x + 3} = x - 3$ bằng
\four{$4$.}{$1$.}{$3$.}{$0$.}

\needspace{4\baselineskip}\item Tích tất cả các nghiệm thực của phương trình $\sqrt{3x^2 - 6x + 1} = \sqrt{2x^2 - 3x - 1}$ bằng
\four{$1$.}{$2$.}{$-2$.}{$3$.}

\end{enumerate}

\newpage

\chapterbanner{PHẦN II. CÂU TRẢ LỜI NGẮN (30 CÂU)}{(Học sinh tính toán, trình bày lời giải chi tiết vào ô li bên dưới và điền kết quả vào ô Đáp số)}
\parttitle{Chủ đề 1. Hàm số và đồ thị (Bài 15)}

\begin{enumerate}

\item Tập xác định của hàm số $y = \sqrt{6 - x} + \dfrac{2x}{\sqrt{x + 2}}$ chứa bao nhiêu số nguyên?\answerbox
\oli{4.5cm}

\item Có bao nhiêu giá trị nguyên của tham số $m \in [-5; 10]$ để hàm số $y = \sqrt{x - m + 1} + \dfrac{1}{\sqrt{2m + 5 - x}}$ xác định trên nửa khoảng $[2; 6)$?\answerbox
\oli{4.5cm}

\item Cho hàm số $f(x) = \begin{cases} 3x - 5 & \text{khi } x \ge 3 \\ \sqrt{7 - x} & \text{khi } -2 \le x < 3 \\ x^2 + 4 & \text{khi } x < -2 \end{cases}$. Tính giá trị biểu thức $T = f(5) + f(-2) + f(-3)$.\answerbox
\oli{4.5cm}

\end{enumerate}

\newpage

\begin{enumerate}[resume*]

\item \begin{minipage}[t]{0.56\linewidth}
Cho hàm số $y = f(x)$ xác định trên đoạn $[-4; 5]$ có đồ thị là đường gấp khúc như hình vẽ bên. Gọi $M, m$ lần lượt là giá trị lớn nhất, nhỏ nhất của hàm số trên đoạn $[-4; 5]$, và $(a; b)$ là khoảng mà hàm số nghịch biến. Tính giá trị của $M + m + a + b$.\answerbox
\end{minipage}\hfill
\begin{minipage}[t]{0.41\linewidth}
\vspace{-6pt}\centering
\begin{tikzpicture}[scale=0.45, >=Stealth]
  \draw[step=1, color=gray!22, line width=0.2pt] (-4.5,-1.5) grid (5.5,5.5);
  \draw[->, line width=0.6pt] (-4.8,0) -- (5.8,0) node[below, font=\scriptsize] {$x$};
  \draw[->, line width=0.6pt] (0,-1.8) -- (0,5.8) node[right, font=\scriptsize] {$y$};
  \node[above right, font=\scriptsize] at (0,0) {$O$};
  \draw[dashed, gray!70, line width=0.4pt] (-4,0) -- (-4,-1) -- (0,-1);
  \draw[dashed, gray!70, line width=0.4pt] (-1,0) -- (-1,5) -- (0,5);
  \draw[dashed, gray!70, line width=0.4pt] (3,0) -- (3,1) -- (0,1);
  \draw[dashed, gray!70, line width=0.4pt] (5,0) -- (5,4) -- (0,4);
  \draw[accent, line width=1.25pt] (-4,-1) -- (-1,5) -- (3,1) -- (5,4);
  \foreach \x/\y in {-4/-1, -1/5, 3/1, 5/4} {
    \fill[accent] (\x,\y) circle (2.6pt);
  }
  \node[above, font=\scriptsize] at (-4,0) {$-4$};
  \node[below, font=\scriptsize] at (-1,0) {$-1$};
  \node[below, font=\scriptsize] at (3,0) {$3$};
  \node[below, font=\scriptsize] at (5,0) {$5$};
  \node[right, font=\scriptsize] at (0,-1) {$-1$};
  \node[left, font=\scriptsize] at (0,1) {$1$};
  \node[left, font=\scriptsize] at (0,4) {$4$};
  \node[right, font=\scriptsize] at (0,5) {$5$};
\end{tikzpicture}
\end{minipage}
\oli{5.0cm}

\item Giá cước của một hãng taxi quy định như sau: $1\text{ km}$ đầu tiên (mở cửa) có giá $15\,000$ đồng; từ trên $1\text{ km}$ đến $15\text{ km}$ có giá $14\,000$ đồng/km; từ trên $15\text{ km}$ trở đi có giá $11\,000$ đồng/km. Một hành khách đi quãng đường $20\text{ km}$ phải trả bao nhiêu nghìn đồng?\answerbox
\oli{5.0cm}

\item Có bao nhiêu giá trị nguyên của tham số $m \in [-6; 6]$ để hàm số bậc nhất $y = (m^2 - 9)x + 3m - 1$ nghịch biến trên $\mathbb{R}$?\answerbox
\oli{5.0cm}

\end{enumerate}

\newpage

\parttitle{Chủ đề 2. Hàm số bậc hai và ứng dụng thực tế (Bài 16)}

\begin{enumerate}[resume*]

\item Cho parabol $(P): y = 3x^2 - 12x + 7$ có đỉnh $I(x_I; y_I)$. Tính giá trị biểu thức $P = x_I - 2y_I$.\answerbox
\oli{5.0cm}

\item Parabol $(P): y = ax^2 + bx + c$ ($a \ne 0$) có đỉnh $I(1; -4)$ và đi qua điểm $A(3; 0)$. Tính giá trị của $a + 2b + 3c$.\answerbox
\oli{5.0cm}

\item \begin{minipage}[t]{0.58\linewidth}
Cho hàm số bậc hai $y = ax^2 + bx + c$ ($a \ne 0$) có đồ thị là parabol như hình vẽ bên. Tính giá trị của biểu thức $2a - b + c$.\answerbox
\end{minipage}\hfill
\begin{minipage}[t]{0.39\linewidth}
\vspace{-6pt}\centering
\begin{tikzpicture}[scale=0.48, >=Stealth]
  \draw[step=1, color=gray!22, line width=0.2pt] (-3.5,-1.5) grid (2.5,4.8);
  \draw[->, line width=0.6pt] (-3.8,0) -- (2.8,0) node[below, font=\scriptsize] {$x$};
  \draw[->, line width=0.6pt] (0,-1.8) -- (0,5.0) node[right, font=\scriptsize] {$y$};
  \node[below right, font=\scriptsize] at (0,0) {$O$};
  \draw[dashed, gray!70] (-1,0) -- (-1,4) -- (0,4);
  \draw[dashed, gray!70] (-1,-1.5) -- (-1,4.8);
  \draw[thick, accent, domain=-3.3:1.3, samples=100] plot (\x, {-(\x)^2 - 2*(\x) + 3});
  \fill[accent] (-1,4) circle (2.4pt) node[above left, font=\scriptsize] {$I(-1;4)$};
  \fill[ink] (-3,0) circle (2pt) node[below left, font=\scriptsize] {$-3$};
  \fill[ink] (1,0) circle (2pt) node[below right, font=\scriptsize] {$1$};
  \fill[ink] (0,3) circle (2pt) node[above right, font=\scriptsize] {$3$};
  \node[right, font=\scriptsize] at (0,4) {$4$};
  \node[below, font=\scriptsize] at (-1,0) {$-1$};
\end{tikzpicture}
\end{minipage}
\oli{5.0cm}

\end{enumerate}

\newpage

\begin{enumerate}[resume*]

\item Tìm giá trị của tham số $m$ để đường thẳng $d: y = 3x + m$ cắt parabol $(P): y = x^2 - x + 2$ tại hai điểm phân biệt $A, B$ sao cho trung điểm $I$ của đoạn thẳng $AB$ có tung độ bằng $8$.\answerbox
\oli{5.5cm}

\item Tìm giá trị nguyên lớn nhất của tham số $m$ để hàm số $y = x^2 - 2(m + 1)x + 5$ đồng biến trên khoảng $(3; +\infty)$.\answerbox
\oli{5.5cm}

\item Một cổng chào hình parabol có chiều cao là $9\text{ m}$ và bề rộng chân cổng ở mặt đất là $6\text{ m}$. Hỏi tại độ cao $5\text{ m}$ so với mặt đất, bề rộng của cổng bằng bao nhiêu mét?\answerbox
\oli{5.5cm}

\end{enumerate}

\newpage

\begin{enumerate}[resume*]

\item Một quả bóng được ném xiên từ độ cao $1{,}5\text{ m}$ theo quỹ đạo parabol $h(x) = -0{,}05x^2 + x + 1{,}5$ (m), với $x$ là khoảng cách theo phương ngang (tính theo mét). Độ cao lớn nhất mà quả bóng đạt được bằng bao nhiêu mét?\answerbox
\oli{5.5cm}

\item Một rạp chiếu phim bán vé với giá $100$ nghìn đồng/vé thì mỗi suất chiếu bán được $120$ vé. Qua khảo sát, cứ mỗi lần giảm giá bán $5$ nghìn đồng/vé thì số vé bán được tăng thêm $10$ vé. Rạp nên đặt giá vé bao nhiêu nghìn đồng để doanh thu của mỗi suất chiếu đạt cao nhất?\answerbox
\oli{5.5cm}

\item Một người thợ có $120\text{ m}$ lưới thép muốn rào thành hai chuồng trại hình chữ nhật bằng nhau nằm sát cạnh nhau (chung một vách ngăn ở giữa, và tận dụng một bức tường dài có sẵn làm cạnh sau không cần rào). Tổng diện tích lớn nhất của hai chuồng trại có thể rào được là bao nhiêu $\text{m}^2$?\answerbox
\oli{5.5cm}

\end{enumerate}

\newpage

\parttitle{Chủ đề 3. Dấu của tam thức bậc hai và bất phương trình (Bài 17)}

\begin{enumerate}[resume*]

\item Tính tổng tất cả các nghiệm nguyên của bất phương trình $-x^2 + 3x + 10 \ge 0$.\answerbox
\oli{5.5cm}

\item Hệ bất phương trình $\begin{cases} x^2 - 6x + 5 \le 0 \\ x^2 - 4x + 3 > 0 \end{cases}$ có bao nhiêu nghiệm nguyên?\answerbox
\oli{5.5cm}

\item Có bao nhiêu số nguyên $x \in [-5; 10]$ là nghiệm của bất phương trình $\dfrac{x^2 - 5x + 4}{x^2 - 16} \le 0$?\answerbox
\oli{5.5cm}

\end{enumerate}

\newpage

\begin{enumerate}[resume*]

\item Có bao nhiêu giá trị nguyên của tham số $m$ để bất phương trình $(m - 1)x^2 - 2(m - 1)x + 3 > 0$ nghiệm đúng với mọi $x \in \mathbb{R}$?\answerbox
\oli{5.5cm}

\item Tìm giá trị nguyên lớn nhất của tham số $m$ để bất phương trình $x^2 - 6x + m \le 0$ nghiệm đúng với mọi $x \in [1; 4]$.\answerbox
\oli{5.5cm}

\item Có bao nhiêu giá trị nguyên của tham số $m \in [-10; 10]$ để bất phương trình $\dfrac{x^2 - 2x + m}{x^2 + 2x + 5} > 0$ nghiệm đúng với mọi $x \in \mathbb{R}$?\answerbox
\oli{5.5cm}

\end{enumerate}

\newpage

\begin{enumerate}[resume*]

\item Một công ty sản xuất $x$ chiếc xe đạp điện mỗi tháng, biết tổng doanh thu là $R(x) = -0{,}5x^2 + 60x$ và tổng chi phí là $C(x) = 20x + 600$ (đơn vị: triệu đồng). Để công ty có lợi nhuận hàng tháng đạt ít nhất $150$ triệu đồng thì số lượng xe $x$ phải nằm trong đoạn $[a; b]$. Tính $b - a$.\answerbox
\oli{5.0cm}

\item Một tấm biển quảng cáo hình chữ nhật có chu vi là $28\text{ m}$. Để diện tích tấm biển không nhỏ hơn $40\text{ m}^2$ và chiều dài lớn hơn chiều rộng, chiều rộng $x$ (m) phải thuộc nửa khoảng $[a; b)$. Tính giá trị của $a + b$.\answerbox
\oli{5.0cm}

\end{enumerate}

\parttitle{Chủ đề 4. Phương trình quy về phương trình bậc hai (Bài 18)}

\begin{enumerate}[resume*]

\item Tìm tổng các nghiệm thực của phương trình $\sqrt{3x^2 - 6x + 4} = \sqrt{2x^2 - x - 2}$.\answerbox
\oli{5.0cm}

\end{enumerate}

\newpage

\begin{enumerate}[resume*]

\item Tìm nghiệm lớn nhất của phương trình $\sqrt{2x^2 + 3x - 5} = \sqrt{x^2 + 5x + 10}$.\answerbox
\oli{5.5cm}

\item Tìm nghiệm thực của phương trình $\sqrt{2x^2 - 7x - 11} = x - 3$.\answerbox
\oli{5.5cm}

\item Tính tổng bình phương tất cả các nghiệm thực của phương trình $\sqrt{5x^2 - 13x + 10} = 2x - 2$.\answerbox
\oli{5.5cm}

\end{enumerate}

\newpage

\begin{enumerate}[resume*]

\item Tính tổng tất cả các nghiệm thực của phương trình $x^2 - 4x - 5\sqrt{x^2 - 4x + 5} + 9 = 0$.\answerbox
\oli{5.5cm}

\item Có bao nhiêu giá trị nguyên của tham số $m \in [-8; 8]$ để phương trình $\sqrt{x^2 - 4x + 13} = x + m$ có nghiệm thực?\answerbox
\oli{5.5cm}

\item \begin{minipage}[t]{0.56\linewidth}
Một hòn đảo $A$ cách bờ biển thẳng một khoảng $AH = 4\text{ km}$. Một trạm phát điện đặt tại điểm $B$ trên bờ biển cách chân đường vuông góc $H$ một đoạn $HB = 12\text{ km}$ (như hình vẽ bên). Người ta kéo dây cáp từ $B$ dọc bờ biển đến điểm $M$ nằm giữa $B$ và $H$, rồi kéo cáp ngầm từ $M$ ra đảo $A$. Biết tổng chiều dài dây cáp sử dụng là $14\text{ km}$. Hỏi điểm $M$ cách $H$ bao nhiêu kilômét?\answerbox
\end{minipage}\hfill
\begin{minipage}[t]{0.41\linewidth}
\vspace{-6pt}\centering
\begin{tikzpicture}[scale=0.38, >=Stealth]
  \fill[cyan!8] (-1,0) rectangle (13,5.2);
  \draw[thick, ink] (-1,0) -- (13,0);
  \coordinate (H) at (0,0);
  \coordinate (A) at (0,4);
  \coordinate (M) at (3,0);
  \coordinate (B) at (12,0);
  \draw[dashed, gray!80] (A) -- (H);
  \draw[line width=1.3pt, accent] (B) -- (M) -- (A);
  \fill[ink] (A) circle (3pt) node[above, font=\scriptsize] {Đảo $A$};
  \fill[ink] (H) circle (2.5pt) node[below, font=\scriptsize] {$H$};
  \fill[accent] (M) circle (2.8pt) node[below, font=\scriptsize] {$M$};
  \fill[ink] (B) circle (2.8pt) node[below, font=\scriptsize] {Trạm $B$};
  \node[left, font=\scriptsize] at (0,2) {$4\text{ km}$};
  \node[above, font=\scriptsize, accent] at (1.5,0) {$x$};
  \draw[<->, gray!80!black] (0,-1.1) -- (12,-1.1) node[midway, below, font=\scriptsize] {$12\text{ km}$};
  \draw (0.5,0) -- (0.5,0.5) -- (0,0.5);
\end{tikzpicture}
\end{minipage}
\oli{5.0cm}

\end{enumerate}

\newpage
\pagestyle{dapan}
\setcounter{page}{1}

\begin{center}
  {\Large\bfseries\color{ink} BẢNG ĐÁP ÁN TỔNG HỢP (50 CÂU)}\\[2pt]
  {\small\color{gray!75!black} Đề cương Ôn tập Chương VI: Hàm số, Đồ thị và Ứng dụng — Toán 10 KNTT}
\end{center}
\vspace{2pt}

\noindent\textbf{PHẦN I. Câu trắc nghiệm nhiều phương án lựa chọn (20 câu)}
\vspace{1.5pt}
\begin{center}
\renewcommand{\arraystretch}{1.15}
\begin{tabular}{|C{0.95cm}|C{0.95cm}|C{0.95cm}|C{0.95cm}|C{0.95cm}|C{0.95cm}|C{0.95cm}|C{0.95cm}|C{0.95cm}|C{0.95cm}|}
\hline
\textbf{1} & \textbf{2} & \textbf{3} & \textbf{4} & \textbf{5} & \textbf{6} & \textbf{7} & \textbf{8} & \textbf{9} & \textbf{10} \\ \hline
\textbf{C} & \textbf{A} & \textbf{D} & \textbf{B} & \textbf{D} & \textbf{B} & \textbf{A} & \textbf{C} & \textbf{B} & \textbf{D} \\ \hline
\hline
\textbf{11} & \textbf{12} & \textbf{13} & \textbf{14} & \textbf{15} & \textbf{16} & \textbf{17} & \textbf{18} & \textbf{19} & \textbf{20} \\ \hline
\textbf{A} & \textbf{C} & \textbf{D} & \textbf{A} & \textbf{C} & \textbf{B} & \textbf{A} & \textbf{D} & \textbf{C} & \textbf{B} \\ \hline
\end{tabular}
\end{center}

\vspace{3pt}
\noindent\textbf{PHẦN II. Câu trắc nghiệm trả lời ngắn (30 câu)}
\vspace{1.5pt}
\begin{center}
\renewcommand{\arraystretch}{1.15}
\begin{tabular}{|C{0.8cm}|C{1.4cm}||C{0.8cm}|C{1.4cm}||C{0.8cm}|C{1.4cm}||C{0.8cm}|C{1.4cm}||C{0.8cm}|C{1.4cm}|}
\hline
\textbf{Câu} & \textbf{Đáp số} & \textbf{Câu} & \textbf{Đáp số} & \textbf{Câu} & \textbf{Đáp số} & \textbf{Câu} & \textbf{Đáp số} & \textbf{Câu} & \textbf{Đáp số} \\ \hline
\textbf{1}  & $8$      & \textbf{7}  & $12$     & \textbf{13} & $6{,}5$  & \textbf{19} & $3$   & \textbf{25} & $5$  \\ \hline
\textbf{2}  & $3$      & \textbf{8}  & $-12$    & \textbf{14} & $80$     & \textbf{20} & $5$   & \textbf{26} & $5$  \\ \hline
\textbf{3}  & $26$     & \textbf{9}  & $3$      & \textbf{15} & $1200$   & \textbf{21} & $9$   & \textbf{27} & $13$ \\ \hline
\textbf{4}  & $6$      & \textbf{10} & $2$      & \textbf{16} & $12$     & \textbf{22} & $20$  & \textbf{28} & $6$  \\ \hline
\textbf{5}  & $266$    & \textbf{11} & $2$      & \textbf{17} & $2$      & \textbf{23} & $11$  & \textbf{29} & $10$ \\ \hline
\textbf{6}  & $5$      & \textbf{12} & $4$      & \textbf{18} & $5$      & \textbf{24} & $5$   & \textbf{30} & $3$  \\ \hline
\end{tabular}
\end{center}

\vspace{4pt}
\begin{center}
  {\large\bfseries\color{ink} HƯỚNG DẪN GIẢI CHI TIẾT}
\end{center}
\vspace{1pt}

\begin{multicols}{2}
\raggedcolumns
\footnotesize

\textbf{\color{accent}PHẦN I. TRẮC NGHIỆM (20 CÂU)}

\textbf{Câu 1.} ĐK: $x + 2 \ge 0 \Leftrightarrow x \ge -2$ và $x^2 - 1 \ne 0 \Leftrightarrow x \ne \pm 1$. Chọn \textbf{C}.

\textbf{Câu 2.} Thay $x = 1 \Rightarrow y = \dfrac{2(1)-1}{\sqrt{1+3}} = \dfrac{1}{2}$. Chọn \textbf{A}.

\textbf{Câu 3.} $f(4) = 3(4) - 2 = 10$; $f(-1) = (-1)^2 + 1 = 2 \Rightarrow P = 12$. Chọn \textbf{D}.

\textbf{Câu 4.} Từ $B(-1; 4)$ đến $C(2; 0)$, khi $x$ tăng thì $y$ giảm nên hàm nghịch biến trên $(-1; 2)$. Chọn \textbf{B}.

\textbf{Câu 5.} Với $x \in [-3; 3]$, $0 \le \sqrt{9 - x^2} \le 3 \Rightarrow 1 \le 4 - \sqrt{9 - x^2} \le 4$. Chọn \textbf{D}.

\textbf{Câu 6.} $x_I = -\dfrac{8}{2(-2)} = 2$, $y_I = -2(4) + 16 - 3 = 5 \Rightarrow I(2; 5)$. Chọn \textbf{B}.

\textbf{Câu 7.} $a = 3 > 0$, trục đối xứng $x = 1 \Rightarrow$ đồng biến trên $(1; +\infty)$. Chọn \textbf{A}.

\textbf{Câu 8.} $x^2 - 5x + 6 = 0 \Leftrightarrow x = 2$ hoặc $x = 3 \Rightarrow AB = |3 - 2| = 1$. Chọn \textbf{C}.

\textbf{Câu 9.} Bề lõm lên $\Rightarrow a > 0$; $x_I = -\dfrac{b}{2a} = -1 \Rightarrow b = 2a > 0$; cắt $Oy$ tại $-3 \Rightarrow c = -3 < 0$. Chọn \textbf{B}.

\textbf{Câu 10.} Đỉnh $x = 1 \in [0; 4]$. $y(0) = -3, y(1) = -4, y(4) = 5 \Rightarrow \max = 5$. Chọn \textbf{D}.

\textbf{Câu 11.} $-\dfrac{b}{2} = 3 \Rightarrow b = -6$; qua $M(1; 0) \Rightarrow 1 - 6 + c = 0 \Rightarrow c = 5 \Rightarrow b + c = -1$. Chọn \textbf{A}.

\textbf{Câu 12.} $x^2 - 4x + 5 = (x - 2)^2 + 1 > 0\ \forall x \in \mathbb{R}$. Chọn \textbf{C}.

\textbf{Câu 13.} $-2x^2 + 5x - 2 > 0$ có $a = -2 < 0$ và hai nghiệm $\dfrac{1}{2}; 2 \Rightarrow x \in \left(\dfrac{1}{2}; 2\right)$. Chọn \textbf{D}.

\textbf{Câu 14.} $3x^2 - 7x + 2 = 0 \Leftrightarrow x = \dfrac{1}{3}$ hoặc $x = 2 \Rightarrow S = \left(-\infty; \dfrac{1}{3}\right] \cup [2; +\infty)$. Chọn \textbf{A}.

\textbf{Câu 15.} $x^2 - 4x - 5 < 0 \Leftrightarrow -1 < x < 5 \Rightarrow x \in \{0; 1; 2; 3; 4\}$ ($5$ nghiệm). Chọn \textbf{C}.

\textbf{Câu 16.} $a = 1 > 0$, yêu cầu $\Delta' = m^2 - m - 2 < 0 \Leftrightarrow -1 < m < 2$. Chọn \textbf{B}.

\textbf{Câu 17.} Bình phương: $x^2 + 3x - 2 = 1 + x \Leftrightarrow x^2 + 2x - 3 = 0 \Leftrightarrow x = 1$ hoặc $x = -3$. Thử lại $x = -3$ làm $1 + x = -2 < 0$ (loại), nhận $x = 1$. Chọn \textbf{A}.

\textbf{Câu 18.} ĐK $x \ge -1$. Bình phương: $2x^2 - 4x + 9 = x^2 + 2x + 1 \Leftrightarrow x^2 - 6x + 8 = 0 \Leftrightarrow x = 2; 4$ ($2$ nghiệm). Chọn \textbf{D}.

\textbf{Câu 19.} ĐK $x \ge 3$. Bình phương: $x^2 - 4x + 3 = x^2 - 6x + 9 \Leftrightarrow 2x = 6 \Leftrightarrow x = 3$ (nhận). Chọn \textbf{C}.

\textbf{Câu 20.} Bình phương: $x^2 - 3x + 2 = 0 \Leftrightarrow x = 1$ hoặc $x = 2$. Thử $x = 1$ cho $3 - 6 + 1 = -2 < 0$ (loại); thử $x = 2$ nhận. Vậy có duy nhất $x = 2$. Chọn \textbf{B}.

\columnbreak

\textbf{\color{accent}PHẦN II. TRẢ LỜI NGẮN (30 CÂU)}

\textbf{Câu 1.} ĐK: $-2 < x \le 6 \Rightarrow x \in \{-1; 0; 1; 2; 3; 4; 5; 6\}$ ($8$ số).

\textbf{Câu 2.} ĐK: $x \ge m - 1$ và $x < 2m + 5$. Để xác định trên $[2; 6)$: $m - 1 \le 2$ và $2m + 5 \ge 6 \Leftrightarrow 0{,}5 \le m \le 3 \Rightarrow m \in \{1; 2; 3\}$ ($3$ giá trị).

\textbf{Câu 3.} $f(5) = 10$, $f(-2) = \sqrt{9} = 3$, $f(-3) = 9 + 4 = 13 \Rightarrow T = 26$.

\textbf{Câu 4.} Từ đồ thị: $M = 5, m = -1$, nghịch biến trên $(-1; 3) \Rightarrow a = -1, b = 3 \Rightarrow 5 - 1 - 1 + 3 = 6$.

\textbf{Câu 5.} $15 + 14 \times 14 + 5 \times 11 = 15 + 196 + 55 = 266$ (nghìn đồng).

\textbf{Câu 6.} Hàm bậc nhất nghịch biến khi $m^2 - 9 < 0 \Leftrightarrow -3 < m < 3 \Rightarrow \{-2; -1; 0; 1; 2\}$ ($5$ giá trị).

\textbf{Câu 7.} $x_I = 2, y_I = 3(4) - 24 + 7 = -5 \Rightarrow P = 2 - 2(-5) = 12$.

\textbf{Câu 8.} $(P): y = a(x - 1)^2 - 4$. Qua $A(3; 0) \Rightarrow 4a - 4 = 0 \Rightarrow a = 1 \Rightarrow y = x^2 - 2x - 3 \Rightarrow a + 2b + 3c = 1 - 4 - 9 = -12$.

\textbf{Câu 9.} Đỉnh $I(-1; 4)$, qua $(0; 3) \Rightarrow y = -(x + 1)^2 + 4 = -x^2 - 2x + 3 \Rightarrow 2(-1) - (-2) + 3 = 3$.

\textbf{Câu 10.} $x^2 - 4x + 2 - m = 0 \Rightarrow x_I = 2$. $I \in d \Rightarrow y_I = 6 + m = 8 \Rightarrow m = 2$ ($\Delta' = 4 > 0$, thỏa).

\textbf{Câu 11.} $a = 1 > 0$, đồng biến trên $(3; +\infty) \Leftrightarrow x_I = m + 1 \le 3 \Leftrightarrow m \le 2 \Rightarrow m_{\max} = 2$.

\textbf{Câu 12.} Chọn trục tọa độ với đỉnh $(0; 9)$, qua $(\pm 3; 0) \Rightarrow y = -x^2 + 9$. Với $y = 5 \Rightarrow x^2 = 4 \Rightarrow x = \pm 2 \Rightarrow$ bề rộng $4\text{ m}$.

\textbf{Câu 13.} Đỉnh $x = -\dfrac{1}{2(-0{,}05)} = 10 \Rightarrow h_{\max} = h(10) = -5 + 10 + 1{,}5 = 6{,}5\text{ m}$.

\textbf{Câu 14.} Gọi số lần giảm $5$ nghìn là $x$: $R(x) = (100 - 5x)(120 + 10x) = -50x^2 + 400x + 12000$. Đạt cực đại tại $x = 4 \Rightarrow$ giá vé $100 - 20 = 80$ nghìn đồng.

\textbf{Câu 15.} Gọi chiều rộng $x \Rightarrow$ dài $120 - 3x$. $S(x) = -3x^2 + 120x$, đạt cực đại tại $x = 20 \Rightarrow S_{\max} = 20 \times 60 = 1200\text{ m}^2$.

\textbf{Câu 16.} $-2 \le x \le 5 \Rightarrow$ tổng $(-2) + (-1) + \dots + 5 = 12$.

\textbf{Câu 17.} $[1; 5] \cap ((-\infty; 1) \cup (3; +\infty)) = (3; 5] \Rightarrow x \in \{4; 5\}$ ($2$ nghiệm).

\textbf{Câu 18.} ĐK $x \ne \pm 4$. Rút gọn $\dfrac{x - 1}{x + 4} \le 0 \Leftrightarrow -4 < x \le 1 \Rightarrow \{-3; -2; -1; 0; 1\}$ ($5$ số).

\textbf{Câu 19.} $m = 1 \Rightarrow 3 > 0$ (đúng $\forall x$). $m \ne 1$: $m - 1 > 0$ và $\Delta' = (m-1)(m-4) < 0 \Leftrightarrow 1 < m < 4$. Vậy $m \in \{1; 2; 3\}$ ($3$ giá trị).

\textbf{Câu 20.} $m \le -x^2 + 6x\ \forall x \in [1; 4]$. Trên $[1; 4]$, $\min(-x^2 + 6x) = 5$ (tại $x=1$) $\Rightarrow m \le 5 \Rightarrow 5$.

\textbf{Câu 21.} Mẫu $x^2 + 2x + 5 > 0\ \forall x$. Tử $> 0\ \forall x \Leftrightarrow 1 - m < 0 \Leftrightarrow m > 1 \Rightarrow \{2;\dots;10\}$ ($9$ giá trị).

\textbf{Câu 22.} $L(x) = -0{,}5x^2 + 40x - 600 \ge 150 \Leftrightarrow x^2 - 80x + 1500 \le 0 \Leftrightarrow 30 \le x \le 50 \Rightarrow 50 - 30 = 20$.

\textbf{Câu 23.} $x(14 - x) \ge 40 \Leftrightarrow 4 \le x \le 10$. Điều kiện rộng $<$ dài $\Leftrightarrow x < 7 \Rightarrow [4; 7) \Rightarrow 4 + 7 = 11$.

\textbf{Câu 24.} $x^2 - 5x + 6 = 0 \Leftrightarrow x = 2; 3$ (cả hai thỏa ĐK) $\Rightarrow 2 + 3 = 5$.

\textbf{Câu 25.} $x^2 - 2x - 15 = 0 \Leftrightarrow x = 5$ hoặc $x = -3$ (cả hai thỏa ĐK) $\Rightarrow$ lớn nhất là $5$.

\textbf{Câu 26.} ĐK $x \ge 3$. Bình phương: $x^2 - x - 20 = 0 \Leftrightarrow x = 5$ (nhận) hoặc $x = -4$ (loại).

\textbf{Câu 27.} ĐK $x \ge 1$. Bình phương: $x^2 - 5x + 6 = 0 \Leftrightarrow x = 2; 3$ (nhận cả hai) $\Rightarrow 2^2 + 3^2 = 13$.

\textbf{Câu 28.} Đặt $t = \sqrt{x^2 - 4x + 5} \ge 1 \Rightarrow t^2 - 5t + 4 = 0 \Leftrightarrow t = 1$ ($x=2$) hoặc $t = 4$ ($x^2 - 4x - 11 = 0$, tổng 2 nghiệm $= 4$). Tổng tất cả các nghiệm $= 2 + 4 = 6$.

\textbf{Câu 29.} ĐK $x \ge -m$. Bình phương: $(2m + 4)x = 13 - m^2$. Cần $m \ne -2$ và $x + m = \dfrac{(m+2)^2 + 9}{2(m+2)} \ge 0 \Leftrightarrow m > -2 \Rightarrow \{-1; 0; \dots; 8\}$ ($10$ giá trị).

\textbf{Câu 30.} Đặt $MH = x$ ($0 \le x \le 12$). Ta có $\sqrt{x^2 + 16} + 12 - x = 14 \Leftrightarrow \sqrt{x^2 + 16} = x + 2 \Leftrightarrow 4x = 12 \Leftrightarrow x = 3\text{ km}$.

\end{multicols}
\end{document}
